% GENERATED FILE. Edit canonical lessons, then run npm run generate:books.
% canonical-source-hash: fnv1a64:39456947fc98cb6e
% canonical-lessons: ES-C56-repaso-amigos

\chapter{Review --- Three Endings}
\label{ch:review-friends}
\hlchaptermodality{26}

\begin{chapteropening}
\textbf{By the end of this chapter:} \emph{I can name the three endings, their partners, and which give gender for free.}

Complete ES-C56-repaso-amigos: give the three pairs and distinguish a decoder from a machine.
\end{chapteropening}

\section[Practice]{Review — three endings}

\label{lesson:ES-C56-repaso-amigos}



\begin{quote}
Nothing new here. Say the three endings and their English partners before you look.
\end{quote}

\begin{grammarlens}[title={What you've built}]
\noindent
\begin{tabularx}{\linewidth}{@{}>{\raggedright\arraybackslash}X>{\raggedright\arraybackslash}X>{\raggedright\arraybackslash}X>{\raggedright\arraybackslash}X@{}}
\toprule
\textbf{Spanish} & \textbf{English} & \textbf{gender} & \textbf{roughly how many} \\
\midrule
\textbf{-ción} & -tion & feminine & \textasciitilde{}2,000 \\
\textbf{-dad} / \textbf{-tad} & -ty & feminine & \textasciitilde{}1,200 \\
\textbf{-mente} & -ly & (adverb) & no limit \\
\bottomrule
\end{tabularx}

Three chapters. Between them they reach further than every vocabulary lesson in this book put together, and they will keep reaching for as long as you keep learning adjectives.

That is worth stating plainly, because it changes what you should expect from the rest of the book. \textbf{A Spanish course is not mostly a list of words.} A large part of it is a small number of rules that turn words you already own into words you can use — and you have just been handed three of the biggest.
\end{grammarlens}

\begin{grammarlens}[title={Grammar Lens: two kinds of ending}]
The three are not the same kind of thing, and the difference matters:

\begin{itemize}
\raggedright
  \item \emph{\textbf{-ción}} and \emph{\textbf{-dad}} are \textbf{decoders}. They let you \emph{read} a word you have never met. They have a fixed reach, because there are only so many English words with those endings.
  \item \emph{\textbf{-mente}} is a \textbf{machine}. It lets you \emph{build} a word nobody taught you, from any adjective you hold. Its reach is whatever your adjective vocabulary turns out to be.
\end{itemize}

Both kinds are worth having, and neither is a substitute for the other. A decoder gets you through a page; a machine gets you through a sentence you have to say yourself.
\end{grammarlens}

\subsection*{Your turn}
\begin{itemize}
\raggedright
  \item \emph{Say it:} \textquotedblleft{}nación, ciudad, rápidamente\textquotedblright{} — one of each
\end{itemize}

\emph{Twice through:}

\begin{tcolorbox}[breakable,colback=teal!4,colframe=teal!35!black,title={Before you move on}]
Three endings — name them. (\emph{\textbf{-ción, -dad, -mente}}.) Which two are always feminine? (\emph{\textbf{-ción}} and \emph{\textbf{-dad}}.) Which one has no upper limit? (\emph{\textbf{-mente}}.) Next: reading a sentence nobody taught you.
\end{tcolorbox}
